% Zdroj: PDF priklady-podzim-2025.pdf, fyzická str. 22-23. Značka: specializace UI-SU, UI-ZPJ. Zařazeno: S3.
\section{Evaluace binárního klasifikátoru}
\szzotazka{podzim 2025}{28}{Evaluace binárního klasifikátoru}

\noindent\textbf{Zadání.}\par
Představte si, že vaším úkolem je vybrat nástroj pro analýzu počítačových logů a vyhledávání anomálií
v nich (ty mohou být způsobené například selháním hardware, kybernetickým útokem, apod.). Máte možnost
si vybrat ze dvou nástrojů. Z dostupných marketingových materiálů jste se o prvním nástroji dozvěděli,
že má úspěšnost (accuracy) $99{,}5\,\%$ a umí správně odhalit $90\,\%$ anomálií. U druhého nástroje se uvádí,
že odhalí $81\,\%$ anomálií a jen $10\,\%$ nahlášených anomálií jsou falešně pozitivní.

\medskip
Předpokládejte, že oba nástroje byly vyhodnoceny na stejných datech obsahujících $10\,000$ instancí,
z nichž pouze $1\,\%$ jsou anomálie.

\begin{enumerate}
  \item Spočítejte chybovou matici (matice konfuze, confusion matrix) obou nástrojů a tyto metriky pro
    binární klasifikaci: úspěšnost (accuracy), přesnost (precision), úplnost (recall), F1-skóre.
  \item Popište výhody a nevýhody jednotlivých nástrojů při praktickém nasazení. Který z nich byste
    použili? Zdůvodněte. (\textit{Zde není nutné jediná správná odpověď.})
\end{enumerate}

\medskip
\noindent\textbf{Náčrt řešení.}\par
O datech víme, že obsahují celkem $9\,900$ negativních instancí (nejsou anomálie) a $100$ pozitivních instancí.
\begin{enumerate}
  \item První nástroj správně detekuje $90$ ze $100$ anomálií. Tedy $\mathrm{TP} = 90$, z toho plyne, že
    $\mathrm{FN} = 10$. Tento nástroj také správně určí $9\,950$ z $10\,000$ instancí. Tedy musí správně určit
    $9\,950 - \mathrm{TP} = 9\,860$ negativních instancí. Tedy $\mathrm{TN} = 9\,860$ a $\mathrm{FP} = 40$.
    Recall je $0.9$ (ze zadání), precision je $\mathrm{TP}/(\mathrm{TP}+\mathrm{FP}) = 90/130$ (cca $0.69$).
    F1-skóre je $0.78$, accuracy (ze zadání) $0.995$.

    \medskip
    Druhý nástroj správně detekuje $80$ ze $100$ anomálií.\footnote{V originálním PDF je zde uvedeno
    \uv{$\mathrm{TP}=81$, a $\mathrm{FN}=19$}, dále \uv{$\mathrm{FN}=81/9=9$} a accuracy $0.997$;
    tyto hodnoty jsou v PDF nekonzistentní (recall $0.81$ při $100$ pozitivních dává $\mathrm{TP}=81$,
    ne $80$). Také je v PDF u precision uveden vzorec $\mathrm{TP}/(\mathrm{TP}+\mathrm{FN})$, správně má
    být $\mathrm{TP}/(\mathrm{TP}+\mathrm{FP})$. Zde přepsáno věrně dle PDF až na opravu vzorce precision;
    případné aritmetické nekonzistence v originálu ponechány.}
    Tedy $\mathrm{TP} = 81$, a $\mathrm{FN} = 19$. Ze zadání jen $10\,\%$ nahlášených jsou FP, tedy
    precision je $0.9 = \mathrm{TP}/(\mathrm{TP}+\mathrm{FP})$ a tedy $\mathrm{FN} = 81/9 = 9$. Zbytek instancí
    ($9\,891$) je TN. Recall je tedy $0.81$ (ze zadání), precision je $0.9$ (ze zadání), F1-skóre je $0.85$,
    accuracy $0.997$.
  \item Druhý nástroj je lepší ve všech metrikách (kromě recall) než nástroj první. Jeho výhodou může být to,
    že hlásí mnohem méně falešně pozitivních výsledků, nezaměstnává tedy správce systému zbytečnými hlášeními.
    Na druhou stranu, pokud jsou anomálie kritické, může být lepší použít první nástroj, který jich detekuje
    více za cenu toho, že okolo $30\,\%$ detekcí jsou falešně pozitivní.
\end{enumerate}

\medskip
\noindent\textit{Doplňující vysvětlení (nad rámec oficiálního náčrtu):} Konzistentní hodnoty pro druhý nástroj: recall $0{,}81$ při $100$ anomáliích dává $\mathrm{TP}=81$, $\mathrm{FN}=19$; precision $0{,}9$ dává $\mathrm{FP}=81/9=9$ (v~náčrtu omylem označeno FN) a~$\mathrm{TN}=9\,900-9=9\,891$. Accuracy je $(81+9\,891)/10\,000=0{,}9972$ a~$F_1=\frac{2\cdot0{,}9\cdot0{,}81}{0{,}9+0{,}81}\approx0{,}85$. U~prvního nástroje je $F_1=\frac{2\cdot(90/130)\cdot0{,}9}{90/130+0{,}9}\approx0{,}78$ a~podíl falešných poplachů mezi hlášeními $40/130\approx31\,\%$.
